\begin{proof}
    Define
    \[
    \veps \coloneqq y - x > 0
    \]
    By \Cref{an:thm:archimedean-property},
    \[
    \exists n \in \N, \frac{1}{n} < \veps
    \]
    Define
    \[
    S \coloneqq \setbuilder*{k \in \Z}{\frac{k}{n} \leq x} \subseteq \Z
    \]
    Then
    \[
    \begin{array}{c}
        \floor{x} \leq x < \floor{x} + 1 \\
        \Downarrow \\
        n\floor{x} \leq nx < n\floor{x} + n \\
        \Downarrow \\
        n\floor{x} \in S \\
        \Downarrow \\
        S \neq \vnot
    \end{array}
    \]
    $S$ is bounded above by $nx$. By \Cref{an:thm:bounded-subset-of-Z-in-R-has-extremum},
    \[
    \exists m \in \Z, m = \max S
    \]
    Since $m = \max S$,
    \[
    \begin{array}{c}
        m \in S \\
        \Downarrow \\
        \frac{m}{n} \leq x \\
        \Downarrow \\
        \frac{m + 1}{n} \leq x + \frac{1}{n}
    \end{array}
    \]
    Also,
    \[
    \begin{array}{c}
        m + 1 \notin S \\
        \Downarrow \\
        \frac{m + 1}{n} > x
    \end{array}
    \]
    To sum up,
    \[
    x < \frac{m + 1}{n} \leq x + \frac{1}{n} < x + \veps = x + (y - x) = y
    \]
    Therefore, we can take $p = \frac{m + 1}{n} \in \Q$ to complete the proof.
\end{proof}

\bigskip

\begin{theorem} \label{an:thm:existence-of-nth-roots}
    \[
    \forall x \in \R^{+}, \forall n \in \N, \exists! y \in \R^{+}, y^{n} = x
    \]
    Such $y$ is denoted by
    \begin{itemize}
        \item $y = \sqrt[n]{x}$
        \item $y = x^{\frac{1}{n}}$
    \end{itemize}
\end{theorem}

\begin{proof}
    Define
    \[
    E \coloneqq \setbuilder*{t \in \R^{+}}{t^{n} < x}
    \]
    If $x \geq 1$, then
    \[
    \begin{array}{c}
        \forall m \in \N, 2^{-m} < 1 \leq x \\
        \Downarrow \\
        2^{-1} \in E \\
        \Downarrow \\
        E \neq \vnot
    \end{array}
    \]
    Otherwise,
    \[
    \begin{array}{c}
        x < 1 \\
        \Downarrow \\
        \forall m \in \N, x^{2m} < x \\
        \Downarrow \\
        x^{2} \in E \\
        \Downarrow \\
        E \neq \vnot
    \end{array}
    \]
    $E$ is bounded above by $\max\set*{x, 1}$. By the \emph{least upper bound property},
    \[
    \exists \alpha \in \R^{+}, \alpha = \sup E
    \]
    We will show that $\alpha^{n} = x$. Let $0 < a < b < \infty$ be given. Then
    \[
    b^{n} - a^{n} = (b - a)\sum_{k = 0}^{n - 1} b^{n - 1 - k} a^{k} < (b - a)nb^{n - 1}
    \]
    Assume that $\alpha^{n} < x$. By \Cref{an:thm:Q-is-dense-in-R},
    \[
    \exists h \in \R^{+}, h < \min\set*{1, \frac{x - \alpha^{n}}{n(\alpha + 1)^{n - 1}}}
    \]
    Put $a = \alpha$ and $b = \alpha + h$. Then
    \[
    \begin{array}{c}
        (\alpha + h)^{n} - \alpha^{n} < hn(\alpha + h)^{n - 1} < hn(\alpha + 1)^{n - 1} \\
        \Downarrow \\
        (\alpha + h)^{n} < \alpha^{n} + hn(\alpha + 1)^{n - 1} < \alpha^{n} + (x - \alpha^{n}) = x
    \end{array}
    \]
    Since $\alpha = \sup E$,
    \[
    \begin{array}{c}
        (\alpha + h)^{n} < x \\
        \Downarrow \\
        \alpha + h \in E \\
        \Downarrow \\
        \bot
    \end{array}
    \]
    Assume that $\alpha^{n} > x$. By \Cref{an:thm:Q-is-dense-in-R},
    \[
    \exists h \in \R^{+}, h < \min\set*{\alpha, \frac{\alpha^{n} - x}{n\alpha^{n - 1}}}
    \]
    Put $a = \alpha - h$ and $b = \alpha$. Then
    \[
    \begin{array}{c}
        \alpha^{n} - (\alpha - h)^{n} < hn\alpha^{n - 1} \\
        \Downarrow \\
        (\alpha - h)^{n} > \alpha^{n} - hn\alpha^{n - 1} > \alpha^{n} - (\alpha^{n} - x) = x
    \end{array}
    \]
    Since $\alpha = \sup E$,
    \[
    \begin{array}{c}
        (\alpha - h)^{n} > x \\
        \Downarrow \\
        \forall t \in E, t^{n} < x < (\alpha - h)^{n} \\
        \Downarrow \\
        \forall t \in E, t < \alpha - h \\
        \Downarrow \\
        \alpha - h \text{ is an upper bound of } E \\
        \Downarrow \\
        \bot
    \end{array}
    \]
    Therefore,
    \[
    \alpha^{n} = x
    \]
    This guarantees the \emph{existence} of $y$. Let $0 < \alpha < \beta < \infty$ be given. Then
    \[
    \begin{array}{c}
        \beta^{n} - \alpha^{n} = (\beta - \alpha)\sum_{k = 0}^{n - 1} \beta^{n - 1 - k} \alpha^{k} > 0 \\
        \Downarrow \\
        \beta^{n} > \alpha^{n}
    \end{array}
    \]
    This guarantees the \emph{uniqueness} of $y$.
\end{proof}

\bigskip

\begin{definition} \label{an:def:greedy-base-expansion-of-real-numbers}
    Let $y \in [0, 1)$ and $b \in \N_{2}$ be given. A sequence $\paren*{a_{n}}_{n = 1}^{\infty}$
    is called a \textbf{greedy base $b$ expansion} of $y$ if
    \[
    \forall n \in \N, 0 \leq y - \sum_{k = 1}^{n} \frac{a_{k}}{b^{k}} < \frac{1}{b^{n}}
    \]
    where
    \[
    \forall k \in \N, a_{k} \in S_{b} \coloneqq \set*{0, 1, 2, \dots, b - 1}
    \]
    In particular, if $b = 10$, such a sequence is called a \textbf{greedy decimal expansion} of $y$.
\end{definition}

\bigskip

\begin{theorem} \label{an:thm:existence-and-uniqueness-of-greedy-base-expansion}
    Let $y \in [0, 1)$ and $b \in \N_{2}$ be given.
    Then the greedy base $b$ expansion of $y$ exists and is unique.
\end{theorem}
